\scanpage{149}
\mspage{V.21}
% archive sheet number: 148

5. \uline{Application des th. de Lefschetz local} (SGA 2 $\overline{\underline{\unsure{\mathrm{XIV}|\mathrm{XV}}}}$ 4).

5.0. Cf. les commentaires dans \unsure{3.3|3.2}. pour les hypothèses faites quand nous \unsure{utiliserons} les théorèmes de Lefschetz locaux.

5.1. Soit $f:X\longrightarrow S$ loc. de type fini, $S$ str. loc. \unsure{hensélien}, et soit $x\in X_0$. Soit $m'\in\mathbb{Z}$ et supposons que l'on \unsure{sache} \deleted{(à titre d'hypothèse}
\[
({\unsure{\underline{\Phi}|\underline{\varphi}}}^i_{\unsure{\nu|v}})_{\bar{x}'}=0
\quad\text{pour tt \unsure{gén. stricte} de $x$ dans $X_0$, et $i\leq m'$.}
\tag*{\tikz[baseline=(n.base)]{\node[inner sep=1pt] (n) {5.1.1};\draw (n.center) ellipse (0.62cm and 0.22cm);}}
\]

\deleted{On} Nous \unsure{considérons} cette hypothèse comme une hypothèse de \unsure{récurrence|voisinage}, et nous proposons de donner des conditions moyennant \deleted{\unsure{un}} \unsure{quoi} on conclura
\[
(\underline{\Phi}^{i})_{\bar{x}}=0\quad\text{si}\quad i\leq m.
\tag*{(*)}
\]

Notons que l'on sait déjà que (5.1.1) implique
\[
\deleted{H^i(\bar{V},\Lambda_\nu)\xrightarrow{\sim}H^i(\unsure{\widetilde{X}_{\bar{\eta}}|\widetilde{X}_{\bar{y}}},\Lambda_\nu)\quad\text{si}\quad i<m'}
\tag*{\deleted{(5.1.1)}}
\]
\deleted{\unsure{en vertu de}}
\[
(\underline{\Phi}^i_\nu)_{\bar{x}}=H^i_x(\widetilde{X},\Lambda_\nu)
=\varinjlim_{\unsure{h|k}} H^i_x(X(h),\Lambda_\nu)
\quad\text{pour}\quad i\leq m'.
\tag*{(5.1.2)}
\]

Par suite, \deleted{\unsure{si on}} la relation (*) \deleted{\unsure{équivaut}} \unsure{résultera} de la relation $H^i_x(X(h),\Lambda_\nu)=0$ pour $i\leq m'$, $\forall h$, i.e. \deleted{\unsure{(5.1.3)}} de la relation
\[
\operatorname{prof\,ét}_x(X(h)\geq m'+1\quad\text{pour}\quad\forall\unsure{k|h}.
\tag*{\tikz[baseline=(n.base)]{\node[inner sep=1pt] (n) {5.1.3};\draw (n.center) ellipse (0.62cm and 0.22cm);}}
\]

\unsure{Donc} On \deleted{\illegible{2}} \deleted{\unsure{en}} en conclut aussitôt, par récurrence sur $\dim{\unsure{\mathscr{O}|\mathcal{O}}}_{X_0,x}=r$ :

\uline{\deleted{\unsure{Proposition}}\inserted{Théorème} 5.2.} Soit $Z$ une partie de $X_0$, \inserted{\unsure{n'importe}.} Supposons que : \deleted{pour}

a) pour $\forall$ générisation $x'$ de $x$ dans $X_0$, \deleted{\unsure{distinct de $x$, on}} contenue dans $X_0-Z$, on ait
\[
(\underline{\Phi}^i_\nu)_{\bar{x}'}=0\quad\text{si}\quad i\leq m',
\]
